We characterise the class of district geometries---``tentacle districts''---where
CH provides diagnostic information that PP and Reock do not.

\subsection{Definition of Tentacle Districts}

A \emph{tentacle district} is a non-convex district in which one or more
thin ``arms'' extend from a main body.
Tentacle geometries arise in hand-drawn maps when map-drawers connect distant
population clusters via thin corridors to achieve partisan or racial targeting
objectives~\citep{stephanopoulos2015partisan}.

\subsection{Synthetic Tentacle Example}

We construct a synthetic tentacle district: a unit square $[0,1]^2$ with
four thin arms of length $5$ and width $0.1$ extending in the cardinal directions.

\paragraph{PP computation.}
Area $A = 1 + 4(5)(0.1) = 3.0$.
Perimeter $P \approx 4(1) + 4 \cdot 2(5 + 0.1) \approx 4 + 40.8 = 44.8$.
PP $= 4\pi(3.0)/(44.8)^2 \approx 37.7/2007 \approx 0.019$.
PP correctly flags this as non-compact.

\paragraph{Reock computation.}
The MBC is determined by the tips of opposite arms.
Distance between opposite arm tips $= 2 \times 5 + 1 = 11$ (approximately),
so $r \approx 5.5$ and MBC area $\approx \pi(5.5)^2 \approx 95.0$.
Reock $= 3.0 / 95.0 \approx 0.032$.
Reock also correctly flags this as non-compact.

\paragraph{Alternative tentacle: moderate PP, low CH.}

Now consider a district with a wider main body ($3 \times 3$ square)
and two thick arms ($4 \times 1$) extending east and west.

Area $A = 9 + 2(4)(1) = 17$.
Perimeter $P = 4(3) + 2 \cdot 2(4 + 1) - 2 \cdot 1 \cdot 2 = 12 + 20 - 4 = 28$.
(Subtracting 4 for the two attachment edges.)
PP $= 4\pi(17)/(28)^2 = 68\pi/784 \approx 0.272$.
MBC: the two easternmost and westernmost points are at distance $3 + 2(4) = 11$,
so $r = 5.5$, MBC area $= \pi(5.5)^2 \approx 95.0$,
Reock $= 17/95 \approx 0.179$.

Convex hull: the convex hull of this shape is an $11 \times 3$ rectangle,
area $= 33$.
CH $= 17/33 \approx 0.515$.

So we have PP $\approx 0.272$ (moderate), Reock $\approx 0.179$ (low),
and CH $\approx 0.515$ (low), confirming Claim~1: some tentacle geometries
have non-trivial PP values while CH correctly detects the non-convexity.

\subsection{Empirical Verification}

\begin{table}[ht]
\centering
\caption{Synthetic tentacle district: PP vs.\ Reock vs.\ CH.}
\label{tab:tentacle}
\begin{tabular}{lccc}
\toprule
District type & PP & Reock & CH \\
\midrule
Thin-arm cross ($w = 0.1$, $L = 5$) & 0.019 & 0.032 & 0.091 \\
Wide-arm dumbbell ($w = 1$, $L = 4$, body $3\times3$) & 0.272 & 0.179 & 0.515 \\
Compact rectangle ($3 \times 4$) & 0.589 & 0.382 & 1.000 \\
\bottomrule
\end{tabular}
\end{table}

The wide-arm dumbbell has PP $\approx 0.27$---a value that would not trigger
concern in many litigation contexts (gerrymandered districts often have
PP $< 0.10$)---but CH $= 0.515$ correctly identifies that 49\% of the
convex hull is unoccupied by the district.
This is the primary diagnostic advantage of CH over PP for tentacle detection.
